% !TeX root = ../main.tex
%------------------------------------------------------------------------------
%  Appendix C -- The companion open-access application
%  The book is deliberately tool-neutral. No method in the main text depends on
%  a particular program; this appendix describes the open-access scripts that
%  accompany the book so that its calculations can be reproduced and varied.
%------------------------------------------------------------------------------
\chapter{The Companion Open-Access Application}
\label{app:software}

No method in this book requires a particular program. Every calculation in the
main text is stated so that it can be followed by hand, and the sequences of
Appendix~\ref{app:calculations} exist for exactly that purpose. Software is a
convenience, not a prerequisite, and a reader with a calculator and the
expressions of Appendix~\ref{app:closed-form} can reproduce every numerical
result printed here.

That said, several of the calculations are tedious by hand and instructive when
varied. \companionname\ therefore accompanies the book. It is a set of scripts
and small interactive applications, released openly at

\begin{center}
\companionlink
\end{center}

\noindent so that the worked examples can be reproduced, the parameters varied
and the sensitivity of each result examined directly. They are provided for
teaching and for checking. None of them contains judgement, and each returns an
answer whether or not its assumptions apply to the problem in hand.

\section{Contents and chapter mapping}
\label{sec:companion-contents}

The scripts are written in Python. Those whose names end in \code{\_app} or
\code{\_streamlit} run as browser-based applications with interactive controls;
the remainder run from the command line and print or plot their results. The
dependencies are listed in \code{requirements.txt} in the repository root.

\begin{longtable}{ll}
\toprule
Script & Supports \\
\midrule
\endhead
\code{Mohr-Coulomb.py} & Chapter~\ref{ch:rock-mechanics} \\
\code{Hoek1998.py} & Chapters~\ref{ch:rock-mechanics}, \ref{ch:ccm} \\
\code{Generalised Hoek-Brown.py} & Chapter~\ref{ch:rock-mass-parameters} \\
\code{MC Parameters from HB.py} & Chapters~\ref{ch:rock-mechanics}, \ref{ch:rock-mass-parameters} \\
\code{GRC.py} & Chapter~\ref{ch:ccm} \\
\code{ccm\_ldp\_scc.py} & Chapter~\ref{ch:ccm} \\
\code{CCM Tunnel.py} & Chapter~\ref{ch:ccm} \\
\code{arc-beam.py} & Chapter~\ref{ch:ccm} \\
\code{tbm\_settlement.py} & Chapter~\ref{ch:numerical-modelling} \\
\code{settlement\_app.py} & Chapter~\ref{ch:numerical-modelling} \\
\code{TBM\_Vibration.py} & Chapter~\ref{ch:numerical-modelling} \\
\code{Retaining-Wall.py} & Chapter~\ref{ch:construction} \\
\bottomrule
\end{longtable}

\section{Structural analysis scripts}
\label{sec:structural-scripts}

\code{Stereonet and LEM.py} and \code{Kinematic Slope.py} contain the
stereographic projection routines used in
Chapter~\ref{ch:structural-analysis}. Although both were written for slope
work, the projection, pole plotting, contouring and plane-intersection
functions apply unchanged to underground wedge and block identification; only
the kinematic failure tests at the end are slope-specific. The net
constructions those routines automate are set out for hand working in
Appendix~\ref{app:stereonet}, and the two should agree.

The remaining slope-specific scripts in the repository belong to companion
slope-engineering material and have no application to the methods in this book.

\section{The demonstrator behind the data figures}
\label{sec:demonstrator-figures}

Seven figures in this book are not drawn but computed: two in
Chapter~\ref{ch:digital-twin}, four in Chapter~\ref{ch:movement-monitoring} and
one in Chapter~\ref{ch:delivery}. They come from a research demonstrator of a
twin-tunnel corridor, and their data is synthetic throughout --- a simulated
alignment through simulated ground, a simulated population of three thousand
buildings above it, and a simulated movement field before and during
construction. No project record appears in any of them.

That is a deliberate choice rather than a limitation. Monitoring records of the
kind these figures show are commercially and legally restricted almost
everywhere, and a book that reproduced one could show it only in part and could
not let a reader vary it. A synthetic corridor can be published whole. The
consequence is that these figures illustrate the shape of a method and not the
behaviour of any ground: no settlement, rate, error or proportion printed in
them should be read as characteristic of a real machine, a real geology or a
real building stock. Where a chapter needs a value a reader might use, it is
given in the text and sourced there.

Two cautions attach to reading them. The demonstrator is seeded, so the figures
are stable between runs and the numbers quoted in the captions match what the
export produces; a change to the demonstrator changes the figures, and the
export script that writes them --- \code{tools/export\_twin\_figures.py} in the
book's own repository --- reports the model error on every run so that such a
change is visible rather than silent. The second caution is that the settlement
bands of Figure~\ref{fig:settlement-screening} are magnitudes of predicted
settlement and not damage categories in the sense of
Section~\ref{sec:damage-assessment}: they carry no building stiffness, no
foundation type and no orientation to the trough, and they stand for a
first-stage screen and nothing beyond it.

\section{Relationship to commercial software}
\label{sec:commercial-software}

Commercial packages are in general use for the analyses described in
Chapters~\ref{ch:ccm} and~\ref{ch:numerical-modelling}, and nothing here argues
against them. They automate the arithmetic, provide meshing and plotting, and
handle geometries that closed-form methods cannot reach.

Two cautions apply, and they are the reason this book states its methods
independently of any package. The first is that a program will return a number
for any input, including inputs outside the range in which its formulation is
valid; the check on that is the reader's, and it is easier to make when the
underlying expression is known. The second is that a result which cannot be
reproduced by an independent route has not been verified. Running a hand
calculation from Appendix~\ref{app:calculations}, or one of the scripts above,
against a package result is the cheapest verification available, and
disagreement between the two is informative whichever turns out to be wrong.
